% =============================================================================
%  main.tex – Hauptdokument (nutzt wbh-studienheft.cls)
% =============================================================================
\documentclass{wbh-studienheft}

% -------------------------  Metadaten  --------------------------------------
\modulkuerzel{ING01-L}
\modultitel{Einführungsprojekt für ingenieurwissenschaftliche Studiengänge}
\modulversion{0721N01}
\autorname{Unger, Jan}
\matrikelnr{863514}

% -------------------------  Zusätzliche Pakete  -----------------------------
% (bei Bedarf ergänzen)
% \usepackage{listings}     % Quellcode-Listings
% \usepackage{amsmath}      % erweiterte Mathematik
% \usepackage{siunitx}      % SI-Einheiten

% =============================================================================
\begin{document}

% ===========================  Titelei  ======================================
\maketitlepage
\makecopyrightpage
\tableofcontents
\clearpage

% ===========================  Vorwort  ======================================
\wbhvorwort{%
  Mit dem Einführungsprojekt für Erstsemester lernen Sie Ihre Hochschule,
  Betreuer und Kommilitonen kennen. Sie bekommen einen ersten Einblick in
  das wissenschaftliche Arbeiten.

  Als Vorbereitung auf diese kleine Projektarbeit sollten Sie dieses
  Studienheft gelesen haben. Theoretische Grundlagen aus technischen
  Fachrichtungen müssen Sie keine mitbringen.
}

% =============================================================================
%  Kapitel 1
% =============================================================================
\chapter{Einführung in die Projekt- und Teamarbeit}

\kapiteleinleitung{In diesem Kapitel erhalten Sie eine Einführung in die
  Projekt- und Teamarbeit, lernen grundlegende Mechanismen in der
  Projekt- und Teamarbeit kennen.}

Mit den folgenden Abschnitten werden einige Grundlagen der Projekt- und
Teamarbeit vorgestellt. Die erfolgreiche Umsetzung eines Projektes in
einem Team ist immer an bestimmte Randbedingungen, Regeln und
Vorgehensweisen gekoppelt.

% ---------------------------------------------------------------------------
\section{Projektdefinition}

Entsprechend DIN~69901 beschreibt man ein Projekt sinngemäß als ein
Vorhaben, das im Wesentlichen durch seine Einmaligkeit gekennzeichnet
ist. Diese Einmaligkeit wird vor allem durch die äußeren Randbedingungen
festgelegt, z.\,B.:

\begin{itemize}[nosep]
  \item klare Zielvorgabe
  \item zeitliche, finanzielle und personelle Begrenzungen
  \item definierter Start- und Endtermin
\end{itemize}

% ---------------------------------------------------------------------------
\section{Aufgaben bei der Projektbearbeitung}

Projekte lassen sich in einzelne Phasen aufteilen: Problemanalyse,
Konzepterstellung, Detailplanung, Realisierung, Produkterstellung und
Abnahme.

\begin{figure}[htbp]
  \centering
  % \includegraphics[width=0.6\textwidth]{bilder/projektmodell.pdf}
  \fbox{\parbox{0.6\textwidth}{\centering\vspace{2cm}%
    \sffamily\small Platzhalter: Dynamisches Projektmodell%
    \vspace{2cm}}}
  \caption{Dynamisches Projektmodell}
  \label{fig:projektmodell}
\end{figure}

\Cref{fig:projektmodell} zeigt das dynamische Projektmodell mit den
Faktoren Qualität, Termine und Kosten.

\begin{table}[htbp]
  \centering
  \caption{Phasen einer Präsentation}
  \label{tab:praesentation}
  \small\sffamily
  \begin{tabularx}{\textwidth}{c l X l c}
    \toprule
    \textbf{Nr.} & \textbf{Modul} & \textbf{Inhalte} & \textbf{Technik} & \textbf{Min.} \\
    \midrule
    1 & Begrüßung   & --                                  & verbal    &  2 \\
    2 & Einstieg     & Darstellung der Projektziele        & Flipchart &  5 \\
    3 & Istzustand   & Aufgabenstruktur                    & Pinnwand  &  5 \\
    4 & Controlling  & Abgearbeitete / offene Aktivitäten  & Pinnwand  & 15 \\
    5 & Aussprache   & Gesamtsituation                     & verbal    & 15 \\
    6 & Aktivitäten  & Maßnahmen, Termine, Zuständige      & Flipchart & 15 \\
    7 & Fazit        & Zusammenfassung                     & verbal    &  5 \\
    \bottomrule
  \end{tabularx}
\end{table}

% ---------------------------------------------------------------------------
\section{Projektmanagement und Teamentwicklung}

Der Teambildungsprozess gliedert sich in fünf Phasen:

\begin{enumerate}[nosep]
  \item Kennenlernphase – \textit{Forming}
  \item Konfliktphase – \textit{Storming}
  \item Stabilisierungsphase – \textit{Norming}
  \item Produktivphase – \textit{Performing}
  \item Projektausstieg – \textit{Ending}
\end{enumerate}

% ---------------------------------------------------------------------------
\section{Präsentationstechnik}

\subsection{Das magische Dreieck der Präsentation}

Das stimmige Zusammenwirken der drei Faktoren Ich, Thema und Gruppe
ist ausschlaggebend für eine erfolgreiche Präsentation.

\subsection{Persönliche Wirkung}

Persönliche Wirkung entsteht durch Präsenz. Präsenz entsteht durch die
glaubwürdige Verkörperung von Extrovertiertheit, Wachsamkeit,
Überzeugungskraft und guter Vorbereitung.

% =============================================================================
%  Anhang
% =============================================================================
\appendix
\listoffigures
\listoftables

% =============================================================================
%  Literaturverzeichnis
% =============================================================================
\begin{thebibliography}{99}
  \bibitem{klan2007}
    Klan, J.; Wittmann, M. (2007).
    \textit{Entwicklung einer Lehrerhandreichung zur Robotertechnik}.
    Münster: Westfälische Wilhelms-Universität Münster.

  \bibitem{scholz2011}
    Scholz, M.\,P. (2011).
    \textit{Roboterwesen bauen und programmieren: Ein Einstieg in
    Lego\textsuperscript{\textregistered} MINDSTORMS\textsuperscript{\textregistered} NXT}.
    Auflage~1, Verlagsgruppe Hüthig-Jehle-Rehm.
\end{thebibliography}

\end{document}
